% =====================================================================
%  บทที่ 4  ผลการวิจัย
% =====================================================================
\chapter{ผลการวิจัย}
\label{ch:results}

\section{ผลการเปรียบเทียบผลสัมฤทธิ์ทางการเรียน}

ผลการเปรียบเทียบคะแนนผลสัมฤทธิ์ทางการเรียนหลังเรียนระหว่างกลุ่มทดลองและกลุ่มควบคุมแสดงในตารางที่~\ref{tab:ttest}

\begin{table}[htbp]
  \caption{ผลการเปรียบเทียบคะแนนหลังเรียนระหว่างกลุ่มทดลองและกลุ่มควบคุม}
  \label{tab:ttest}
  \centering
  \begin{tabular}{@{}l c c c c c@{}}
    \toprule
    \textbf{กลุ่ม} & $n$ & $\bar{x}$ & $S.D.$ & $t$ & $p$\\
    \midrule
    กลุ่มทดลอง  & 30 & 24.37 & 2.81 & \multirow{2}{*}{4.12} & \multirow{2}{*}{.000*}\\
    กลุ่มควบคุม & 30 & 21.13 & 3.24 &  & \\
    \bottomrule
    \multicolumn{6}{@{}l}{\small * มีนัยสำคัญทางสถิติที่ระดับ .05}
  \end{tabular}
\end{table}

จากตารางที่~\ref{tab:ttest} พบว่ากลุ่มทดลองมีคะแนนเฉลี่ยหลังเรียนสูงกว่ากลุ่มควบคุมอย่างมีนัยสำคัญทางสถิติที่ระดับ .05 และเมื่อพิจารณาคะแนนรายหัวข้อ พบว่าหัวข้อที่มีความแตกต่างมากที่สุดคือการวนซ้ำ ดังภาพที่~\ref{fig:scores}

\begin{figure}[htbp]
  \centering
  \begin{tikzpicture}
    \begin{axis}[
        ybar, bar width=9pt, width=0.9\linewidth, height=6.5cm,
        ymin=0, ymax=10, ylabel={คะแนนเฉลี่ย},
        symbolic x coords={ตัวแปร, เงื่อนไข, การวนซ้ำ, ฟังก์ชัน, รายการ},
        xtick=data, enlarge x limits=0.15,
        legend style={at={(0.5,-0.18)}, anchor=north, legend columns=-1,
                      draw=none, /tikz/every even column/.append style={column sep=1em}},
        tick label style={font=\small}, label style={font=\small}]
      \addplot[fill=blue!55]   coordinates {(ตัวแปร,8.9) (เงื่อนไข,8.1) (การวนซ้ำ,7.6) (ฟังก์ชัน,7.2) (รายการ,7.0)};
      \addplot[fill=orange!60] coordinates {(ตัวแปร,8.5) (เงื่อนไข,7.4) (การวนซ้ำ,5.9) (ฟังก์ชัน,6.3) (รายการ,6.1)};
      \legend{กลุ่มทดลอง, กลุ่มควบคุม}
    \end{axis}
  \end{tikzpicture}
  \caption{คะแนนเฉลี่ยรายหัวข้อของกลุ่มทดลองและกลุ่มควบคุม}
  \label{fig:scores}
\end{figure}

\section{ผลการประเมินความพึงพอใจ}

ตารางที่~\ref{tab:satisfaction} แสดงตัวอย่างตารางยาวที่ต่อข้ามหน้าได้ (\texttt{longtable}) โดยหัวตารางจะแสดงซ้ำพร้อมคำว่า ``(ต่อ)'' ในหน้าถัดไปโดยอัตโนมัติ

\begin{longtable}{@{}p{0.62\linewidth} c c l@{}}
  \caption{ความพึงพอใจของผู้เรียนที่มีต่อระบบผู้ช่วยสอน}
  \label{tab:satisfaction}\\
  \toprule
  \textbf{รายการประเมิน} & $\bar{x}$ & $S.D.$ & \textbf{ระดับ}\\
  \midrule
  \endfirsthead
  \caption[]{(ต่อ)}\\
  \toprule
  \textbf{รายการประเมิน} & $\bar{x}$ & $S.D.$ & \textbf{ระดับ}\\
  \midrule
  \endhead
  \bottomrule
  \endlastfoot
  1. ระบบตอบคำถามได้ตรงกับเนื้อหาของรายวิชา & 4.60 & 0.50 & มากที่สุด\\
  2. คำอธิบายเข้าใจง่าย & 4.53 & 0.57 & มากที่สุด\\
  3. ระบบช่วยให้แก้ไขข้อผิดพลาดของโค้ดได้เร็วขึ้น & 4.57 & 0.63 & มากที่สุด\\
  4. ระบบตอบสนองได้รวดเร็ว & 4.37 & 0.67 & มาก\\
  5. ระบบใช้งานง่าย & 4.63 & 0.49 & มากที่สุด\\
  \midrule
  \textbf{ภาพรวม} & \textbf{4.52} & \textbf{0.58} & \textbf{มากที่สุด}\\
\end{longtable}
